\documentclass[10pt,letterpaper,final]{report}
\usepackage[margin=.75in]{geometry}
\usepackage[utf8]{inputenc}
\usepackage{amsmath}
\usepackage{amsfonts}
\usepackage{amssymb}
\usepackage{amsthm}
\renewcommand\qedsymbol{$\blacksquare$}
\usepackage{enumerate}
\usepackage{hyperref}
\usepackage{pdfpages}
\usepackage{graphics}
\usepackage{graphicx}
\usepackage{tikz}
\usepackage{tikz-qtree}
\usetikzlibrary{automata,arrows}
\usepackage{mdframed}
\usepackage[
  backend=biber,
  style=numeric-comp,
  sorting=none
]{biblatex}
\addbibresource{references.bib}

\usepackage{minted}
\usemintedstyle{colorful}
\usepackage{xltabular}
\usepackage{pdflscape}


% Based on template from COSC-417 at Towson University, originally authored by Marius Zimand

\begin{document}
\begin{center}
  \fbox{
    \vbox{
      \begin{flushleft}
        Jonathan Llovet \\  % authors' names
        EN.605.202.81 \\  %class
        2026-06-30 \\  % date
        Module 5 - Homework 5 \\ % ID
      \end{flushleft}
      \center{\Large{\textbf{Data Structures - More on Lists}}}
      %\end{mdframed}
    } % end vbox
  } % end fbox
  \vline
\end{center}

Write pseudo-code not Java for problems requiring code. You are responsible for the appropriate level of detail.

The questions in this assignment give you the opportunity to explore a new data structure and to experiment with the hybrid implementation in Q3.

\medskip

\noindent\textbf{Problem 1:}

A deque (pronounced deck) is an ordered set of items from which items may be deleted at either end and into which items may be inserted at either end. Call the two ends left and right. This is an access-restricted structure since no insertions or deletions can happen other than at the ends. Implement the deque as a doubly-linked list (not circular, no header). Write InsertLeft and DeleteRight.

\medskip
\noindent\textbf{Answer:}
\medskip

\pagebreak

\noindent\textbf{Problem 2:}

Implement a deque from problem 1 as a doubly-linked circular list with a header. Write InsertRight and DeleteLeft.


\medskip
\noindent\textbf{Answer:}
\medskip

\pagebreak

\noindent\textbf{Problem 3:}

Write a set of routines for implementing several stacks and queues within a single array. Hint: Look at the lecture material on the hybrid implementation.


\medskip
\noindent\textbf{Answer:}
\medskip

\pagebreak

\section*{Source Code}
The full source code, including LaTeX, tooling, other artifacts, and git history is available on my Github repository for the class here:
\begin{center}
  \url{https://github.com/jllovet/jhu-en-605-202-su26-data-structures}
\end{center}

\printbibliography

\end{document}
